\chapter{Ideas and Concepts}

\section{Overall Pipeline Concept}
Bridging the data gap identified in (REFERENCE SECTION 1.2) requires an end-to-end workflow that connects 
synthetic data generation to detector training and evaluation under a unified benchmark. The proposed pipeline 
takes a scarce base dataset of real images as its starting point and produces a set of dataset variants through 
progressively synthetic augmentation strategies, ranging from classical geometric and photometric transforms to fully
 generative image synthesis.
 Each variant is used to train a \gls{yolo} detector under identical hyperparameters, comparing its performance systematically. 
Synthetic data quality is therefore not assessed in isolation but through its downstream effect on \gls{map@50}, precision, recall, and false negative rate. 

Two diffusion approaches are adopted, described in detail in [Sections 3.2] and [3.3] respectively. One generates new images by interpolating between a the real 
base distribution and a text-conditioned synthetic one, while the other composites target object instances into new background scenes via inpainting. Both are 
trained on the same scarce base dataset, and their output will be used to create the dataset variants.

While advances in synthetic data generation have produced promising alternatives, the pipeline builds on diffusion-based image synthesis as its generative foundation. 
Three-dimensional rendering and game-engine based synthesis have demonstrated strong results in controlled indoor settings, where scene geometry, lighting, 
and object placement can be easily parameterized. The Ascento Guard operates almost exclusively outdoors and regularly covers missions of five kilometers and more, driving 
through environments whose variability in terrain, vegetation, wheather conditions, and background clutter cannot be captured by any fixed scene graph. Modeling this exhaustively 
per customer size would require a manual authoring effort that defeats the purpose of a generalizable pipeline. Diffusion-based synthesis sidesteps this problem by operating 
directly in image space, inheriting the statistical properties of the real base dataset rather then attempting to reconstruct the deployment environement from geometric primitives. 

\section{Approach A: Diffusion Curriculum}
Standard text-guided diffusion synthesis lacks control over how closely generated images resemble the originals. When real training data is scarce or low-quality, text-only synthesis 
introduces a distribution gap that hurts downstream classifier performance. \gls{discl}\cite{diffusion-curriculum} addresses this by using image guidance withing the reverse diffusion process to 
produce a continuous spectrum of interpolations between a real image and a purely text-generated one. 
The reverse diffusion process follows the standard \gls{ddpm} formulation. 
At each timestep t the noise estimate is computed via classifier-free guidance:

\begin{equation}
    \label{eq:cfg-noise}
    \hat{\varepsilon}_t = (1 + w) \varepsilon_\theta(z_t, t \mid c) - w \varepsilon_\theta(z_t, t)
\end{equation}
where $c$ s the text prompt injected into the \gls{u-net} vis cross-attention.
$w$ controls the strength of that conditioning, and $\epsilon_\theta(z_t, t)$ is the unconditional prediction. The denoising step then produces $z_{t-1}$. 

% CHECK THIS AGAIN
\begin{equation}
    \label{eq:denoising-step}
    z_{t-1}= (\frac{1}{\sqrt{\alpha_t}})(z_t - (\frac{\beta_t} {\sqrt{1 - \tilde{\alpha}_t}}) \hat{\epsilon}_t) + \sqrt{\beta_t \epsilon}
\end{equation}
The variance schedule $\tilde{\alpha}_t, \alpha_t, \beta_T$  is fixed across all timesteps using \gls{ddim} and monotonically decreasing $\tilde{\alpha}_t$ 
is what makes the image guidance work. The image guidance leve $\lambda \in [0,1]$ control where the reverse diffusion process begins. Rather than starting from \gls{isotropic Gaussian noise} $x_T$ 
the latent of the real image $z_{real}$ is perturbed to timestep $t(\lambda)=\lfloor (1 -\lambda)\rfloor$:
\begin{equation}
    \label{eq:image-guidance}
    z_{t(\lambda)} = \sqrt{\tilde{\alpha}_{t(\lambda)}} + \sqrt{(1 - \tilde{\alpha}_{t(\lambda)})}\cdot \epsilon, \epsilon \sim \mathcal{N}(0, 1)
\end{equation}

From this starting point, standard classifier-free guidance denoising proceeds conditioned on a text prompt $c$. A large $\lambda$ means $t(\lambda)$ is small; little noise is injected and $z_{real}$ dominates, 
so the output closely resembles the original. A small $\lambda$ corresponds to a large $t(\lambda)$ resulting in the real image being heavily corrupted before denoising begins and the output is driven almost 
entirely by the text prompt. This produces a full spectrum of synthetic data, varying jointly in diversity, quality, and proximity to the real distribution. The text prompt $c$, carrying class-level semantic 
information, is injected at every denoising step regardless of $\lambda$. Generated images are then filtered for semantic fidelity using \gls{clipscore} and discarded if found to be below a task specific 
threshold before the training begins. The curriculum then selects which guidance level to expose the model to at each training stage. For long-tail classification where tail classes suffer from data scarcity 
and low diversity, a non-adaptive schedule is used. Training begins with low $\lambda$ synthetic images to exploit their diversity, and $\lambda$ is increased linearly until the data resembles the original 
distribution. For low-quality data the optimal guidance level is not known in advance. An adaptive curriculum instead evaluates at each epoch the improvement in true-class confidence on a held out 
validation subset for each available $\lambda$, then selects the guidance level that produced the greatest progress for the next training stage.
On \gls{imagenet}, \gls{discl} raised tail-class accuracy from $4.4\%$ to $23.64\%$ over the cross-entropy baseline. On \gls{iwildcam} it achieves a 2.7\% \gls{ood} and 2.1\% ID \gls{macrof1} improvement over the strongest diffusion-augmentation baseline. 

\section{Approach B: Multimodel object inpainting}
\gls{mobi}\cite{mobi} extends Paint-by-Example to condition the reverse diffusion process on a projected 3D bounding box alongside the standard \gls{clip} reference token. The full conditioning vector at each denoising step is
\begin{equation}
    \label{eq:mobi-conditioning}
    c = {c_{env}^C, c_{env}^R, c_{ref}, c_{box}^C, c_{box}^R, \tilde{m}^C, \tilde{m}^R}
\end{equation}

EXPLAIN VARIABLES!!!!!!!!!
The model is trained to minimise the noise prediction loss over both camera and \gls{lidar} simultaneously. The eight 3D bounding box corner are projected into image space and augmented with depth from the \gls{lidar} projection to produce a depth-aware camera box. 
Both the camera and LiDAR projections are encoded into conditioning tokens via \gls{fourier embeddings} followed by modality-specific trainable linear layers.
These tokens are injected through a dedicated gated bounding box adapter interleaved with a frozen \gls{pbe} cross-attention layers, using zero-initialized gating as \gls{controlnet}. 
Decoupling the adapter from the \gls{pbe} cross-attention layers keeps the bounding box spatial signale from interfering with the reference appearance already encoded in $c_{ref}$. 
Without this separation, concatenating the box token directly into the \gls{pbe} cross-attention layer would degrade the perceptual alignment.  
To enforce geometric consistency between camera and \gls{lidar} outputs, a cross-modal attention layer is interleaved at each \gls{u-net} block.
Camera features $h^C$ attend to \gls{lidar} $h^R$  features via learnable projections $W_Q^C, W_K^R, W_V^R$.

\begin{equation}
    \label{eq:cross-modal-attention}
    Attention^C = softmax(\frac{Q^C(K^R)^T} {\sqrt{d_{head}})} V^R
\end{equation}
with a residual update $h^C \leftarrow h^C + Gate^C(Attention^C)$ through zero-initialized gate. The network learns implicit alignment facilitated by the shared bounding box projections appearing both in 
$c_{box}^C$ and $c_{box}^R$ 

During training source and destination bounding boxes correspond to the same object at different timestamps, while at inference they can be chosen randomly. Reference are sampled with a \gls{beta(4,1)} 
distribution over normalized temporal distance, biasing towards frames far from the insertion time stamp to prevent the network from learning trivial identity mapping. Empty bounding boxes with 
with black reference images are included in 30\% of the training samples, forcing the model to reconstruct background from context alone. The autoencoders and all \gls{pbe} layers remain frozen: 
only the bounding box encoder, bounding box adapter, and cross-modal attention layers are trained. 
